% TODO make hw stylesheet
\documentclass[10pt]{scrartcl}
\usepackage[a4paper, total={6in, 8in}]{geometry}
\usepackage[usenames,dvipsnames,svgnames]{xcolor}
\usepackage[shortlabels]{enumitem}
\usepackage[framemethod=TikZ]{mdframed}
\usepackage{amsmath,amssymb,amsthm}
\usepackage[colorlinks]{hyperref}
\usepackage{multicol}
\usepackage{tabularx}

\hypersetup{
	colorlinks=true,
	linkcolor=blue,
	filecolor=magenta,      
	urlcolor=magenta,
	pdftitle={MAT 362 HW}
}

\usepackage{microtype}
\usepackage{mathtools}
\usepackage[headsepline]{scrlayer-scrpage}
\usepackage{thmtools}
\usepackage[textsize=footnotesize]{todonotes}

\mdfdefinestyle{mdbluebox}{roundcorner=10pt,innerbottommargin=9pt,
	linecolor=blue,backgroundcolor=TealBlue!5,}
\declaretheoremstyle[headfont=\sffamily\bfseries\color{MidnightBlue},
mdframed={style=mdbluebox},]{thmbluebox}

\mdfdefinestyle{mdbloobox}{frametitlefont=\bfseries,innerbottommargin=8pt,
	nobreak=true,backgroundcolor=TealBlue!5,linecolor=blue,}
\declaretheoremstyle[headfont=\bfseries\color{MidnightBlue},
mdframed={style=mdbloobox},headpunct={\\[3pt]},postheadspace=0pt,]{thmbloobox}

\mdfdefinestyle{mdredbox}{frametitlefont=\bfseries,innerbottommargin=8pt,
	nobreak=true,backgroundcolor=Salmon!5,linecolor=RawSienna,}
\declaretheoremstyle[headfont=\bfseries\color{RawSienna},
mdframed={style=mdredbox},headpunct={\\[3pt]},postheadspace=0pt,]{thmredbox}

\mdfdefinestyle{mdgreenbox}{linecolor=ForestGreen,backgroundcolor=ForestGreen!5,
	linewidth=2pt,rightline=false,leftline=true,topline=false,bottomline=false,}
\declaretheoremstyle[headfont=\bfseries\sffamily\color{ForestGreen!70!black},
mdframed={style=mdgreenbox},headpunct={ --- },]{thmgreenbox}

\mdfdefinestyle{mdorangebox}{linecolor=RedOrange,backgroundcolor=RedOrange!5,
	linewidth=2pt,rightline=false,leftline=true,topline=false,bottomline=false,}
\declaretheoremstyle[headfont=\bfseries\sffamily\color{RedOrange!70!black},
mdframed={style=mdorangebox},headpunct={ --- },]{thmorangebox}

\mdfdefinestyle{mdblackbox}{linecolor=black,backgroundcolor=RedViolet!5!gray!5,
	linewidth=3pt,rightline=false,leftline=true,topline=false,bottomline=false,}
\declaretheoremstyle[mdframed={style=mdblackbox}]{thmblackbox}

\declaretheoremstyle[spaceabove=6pt,spacebelow=6pt]{basehead}

\declaretheorem[style=basehead,name=Problem]{problem}
\declaretheorem[style=basehead,name=Problem,numbered=no]{problem*}
\declaretheorem[style=thmbloobox,name=Setup]{setup} % for config geo
\declaretheorem[style=thmredbox,name=Example,sibling=problem]{example}
\declaretheorem[style=thmredbox,name=$\bigstar$ Problem,sibling=problem]{niceprob}
\declaretheorem[style=thmbluebox,name=Theorem,sibling=problem]{theorem}
\declaretheorem[style=thmbluebox,name=Corollary,sibling=theorem]{corollary}
\declaretheorem[style=thmbluebox,name=Proposition,sibling=theorem]{prop}
\declaretheorem[style=thmbluebox,name=Lemma,numberwithin=problem]{lemma}
\declaretheorem[style=thmbluebox,name=Theorem,numbered=no]{theorem*}
\declaretheorem[style=thmbluebox,name=Lemma,numbered=no]{lemma*}
\declaretheorem[style=thmgreenbox,name=Claim,numberwithin=problem]{claim}
\declaretheorem[style=thmgreenbox,name=Claim,numbered=no]{claim*}
\declaretheorem[style=thmblackbox,name=Remark,numberwithin=problem]{remark}
\declaretheorem[style=thmblackbox,name=Remark,numbered=no]{remark*}
\declaretheorem[style=thmgreenbox,name=Definition,sibling=theorem]{definition}
\declaretheorem[style=thmgreenbox,name=Definition,numbered=no]{definition*}
\declaretheorem[style=thmorangebox,name=Warning,sibling=theorem]{warning}
\declaretheorem[style=thmorangebox,name=Warning,numbered=no]{warning*}
\declaretheorem[style=thmorangebox,name=Note,numbered=no]{note}
\declaretheorem[style=thmblackbox,name=Exercise,sibling=problem]{exercise}
\declaretheorem[style=thmblackbox,name=Exercise,numbered=no]{exercise*}
\declaretheorem[style=thmblackbox,name=Question,sibling=problem]{question}
\declaretheorem[style=thmblackbox,name=Question,numbered=no]{question*}

\addtolength{\textheight}{3.14cm}
\providecommand{\ol}{\overline}
\providecommand{\ul}{\underline}
\providecommand{\eps}{\varepsilon}
\providecommand{\half}{\frac{1}{2}}
\providecommand{\dang}{\measuredangle} %% Directed angle
\providecommand{\CC}{\mathbb C}
\providecommand{\FF}{\mathbb F}
\providecommand{\NN}{\mathbb N}
\providecommand{\QQ}{\mathbb Q}
\providecommand{\RR}{\mathbb R}
\providecommand{\ZZ}{\mathbb Z}
\providecommand{\dg}{^\circ}
\providecommand{\ii}{\item}
\providecommand{\setdiff}{\smallsetminus}
\providecommand{\alert}[1]{{\sffamily\textbf{\textcolor{blue}{#1}}}}
\providecommand{\inv}{^{-1}}
\providecommand{\mb}[1]{\mathbf{#1}}
\providecommand{\dif}{\;d}
\newcommand*{\horzbar}{\rule[.5ex]{2.5ex}{0.5pt}}

\reversemarginpar

\newenvironment{soln}{\begin{proof}[Solution]}{\end{proof}}
\newenvironment{walkthrough}{\noindent \textbf{Walkthrough.}}{\medskip}
\newlist{walk}{enumerate}{3}
\setlist[walk]{label=\bfseries (\alph*)}

\providecommand{\pow}{\operatorname{Pow}}
\providecommand{\vocab}[1]{{\textbf{\textcolor{ForestGreen}{#1}}}}
\providecommand{\lcm}{\operatorname{lcm}}
\providecommand{\abs}[1]{\left|#1\right|}

\usepackage{asymptote}
\begin{asydef}
	defaultpen(fontsize(10pt));
	unitsize(1cm);
	size(8cm); // set a reasonable default
	usepackage("amsmath");
	usepackage("amssymb");
	settings.tex="pdflatex";
	settings.outformat="pdf";
	// Replacement for olympiad+cse5 which is not standard
	import geometry;
	// recalibrate fill and filldraw for conics
	void filldraw(picture pic = currentpicture, conic g, pen fillpen=defaultpen, pen drawpen=defaultpen)
	{ filldraw(pic, (path) g, fillpen, drawpen); }
	void fill(picture pic = currentpicture, conic g, pen p=defaultpen)
	{ filldraw(pic, (path) g, p); }
	// some geometry
	pair foot(pair P, pair A, pair B) { return foot(triangle(A,B,P).VC); }
	pair orthocenter(pair A, pair B, pair C) { return orthocentercenter(A,B,C); }
	pair centroid(pair A, pair B, pair C) { return (A+B+C)/3; }
	// cse5 abbreviations
	path CP(pair P, pair A) { return circle(P, abs(A-P)); }
	path CR(pair P, real r) { return circle(P, r); }
	pair IP(path p, path q) { return intersectionpoints(p,q)[0]; }
	pair OP(path p, path q) { return intersectionpoints(p,q)[1]; }
	path Line(pair A, pair B, real a=0.6, real b=a) { return (a*(A-B)+A)--(b*(B-A)+B); }
	// cse5 more useful functions
	picture CC() {
		picture p=rotate(0)*currentpicture;
		currentpicture.erase();
		return p;
	}
	pair MP(Label s, pair A, pair B = plain.S, pen p = defaultpen) {
		Label L = s;
		L.s = "$"+s.s+"$";
		label(L, A, B, p);
		return A;
	}
	pair Drawing(Label s = "", pair A, pair B = plain.S, pen p = defaultpen) {
		dot(MP(s, A, B, p), p);
		return A;
	}
	path Drawing(path g, pen p = defaultpen, arrowbar ar = None) {
		draw(g, p, ar);
		return g;
	}
\end{asydef}

\usepackage{mathrsfs}

\ihead{\footnotesize MAT 362 HW11}
\ohead{\footnotesize Last compiled \today}

\title{\vspace{-2em}MAT 362 HW11}
\author{\large Aditya Pahuja}
\date{\large Due 05/01}
\begin{document}
\maketitle

\begin{problem*}4.5.1.
	Let $S\subseteq\RR^3$ be a regular, compact, connected, orientable surface
	which is not homeomorphic to a sphere. Prove that there are points on $S$
	where the Gaussian curvature is positive, negative, and zero.
\end{problem*}

\begin{soln}
	We first show that any compact surface has a point of positive curvature.
	Let $p\in S$ be a point of $S$ whose distance from the origin is maximal
	(which exists by compactness). Then $S$ lies entirely inside
	one of the two half-spaces demarcated by the plane $p+T_p(S)$,
	touching this tangent plane at exactly one point
	(notably $p$ is normal to $T_p(S)$).
	This implies that the normal to any normal section at $p$ will have
	the same direction as the vector $p$ (and not the opposite direction),
	so both principal curvatures have the same sign, hence $p$
	is a point at which Gaussian curvature is positive.
	In fact, the Gaussian curvature is positive in a neighborhood of $p$
	because the curvature varies continuously.
	(More rigorously, one could apply a rotation and scale
	so that $p = (0,0,1)$ and $T_p(S)$ is the $xy$-plane, then see that
	$S$ is the graph of a function with maximum $1$ in a neighborhood of $p$.)
	
	Since $S$ is not homeomorphic to a sphere,
	its Euler-Poincar\'e characteristic is nonpositive,
	so the integral of $K$ over $S$ is nonpositive.
	The existence of an open set with positive Gaussian curvature at every point
	implies the existence of an open set with negative values of $K$
	(else the integral would be positive), and the intermediate value theorem
	shows us that there exists a point with curvature zero
	(say, by taking a path in $S$ between a point with positive $K$
	and a point with negative $K$).
\end{soln}

\begin{problem*}4.5.2.
	Let $T$ be a torus of revolution. Describe the image of
	the Gauss map of $T$ and show, without using the Gauss-Bonnet theorem, that
	\[\iint_TK\dif\sigma=0.\]
	Compute the Euler-Poincar\'e characteristic of $T$ and check
	the above result with the Gauss-Bonnet theorem.
\end{problem*}

\begin{soln}
	The image of the Gauss map of $T$ is the entire sphere,
	since we can cover a point
	$(\cos\theta\sin\phi,\sin\theta\sin\phi,\cos\phi)$
	on the sphere by taking the cross section given by
	rotating the generator of the torus in the $xz$-plane by $\theta$
	and then moving an angle of $\phi$ within this cross section.
	
	The characteristic of $T$ is $0$, as shown by the picture in the text.
\end{soln}

\begin{problem*}4.5.3.
	Let $S\subseteq\RR^3$ be a regular compact surface with $K>0$.
	Let $\Gamma\subseteq S$ be a simple closed geodesic in $S$,
	and let $A$ and $B$ be the regions of $S$ which have $\Gamma$
	as a common boundary. Let $N\colon S\to S^2$ be the Gauss map of $S$.
	Prove that $N(A)$ and $N(B)$ have the same area.
\end{problem*}

\begin{soln}
	Since the geodesic curvature of $C$ is $0$ everywhere
	and it is easily checked that the characteristic of $A$ is $1$,
	Gauss-Bonnet gives
	\[\iint_A K\dif\sigma = 2\pi.\]
	Then a change of variables gives
	\[\iint_{N(A)} 1\dif\ol\sigma = \iint_A K\dif\sigma = 2\pi\]
	so $N(A)$ and $N(B)$ both have area $2\pi$
	because the total surface area of $S^2$ is $4\pi$.
\end{soln}

\begin{problem*}4.5.4.
	Compute the Euler-Poincar\'e characteristic of an ellipsoid.	
\end{problem*}

\begin{soln}
	Ellipsoids are diffeomorphic to spheres, so the characteristic of
	an ellipsoid is the same as that of a sphere, namely \framebox{2}.
\end{soln}


\begin{problem*}4.5.5.
	Let $C$ be a parallel of colatitude $\varphi$ on
	an oriented unit sphere $S^2$, and let $w_0$ be
	a unit vector tangent to $C$ at a point $p\in C$.
	Take the parallel transport of $w_0$ along $C$ and show that
	its position, after a complete turn,
	makes an angle $\Delta\varphi = 2\pi(1-\cos\varphi)$
	with the initial position $w_0$. Check that
	\[\lim_{R\to p} \frac{\Delta\varphi}{A} = 1 = \text{curvature of }S^2\]
	where $A$ is the area of the region $R$ of $S^2$ bounded by $C$.
\end{problem*}

\begin{soln}
	We have that the radius of $C$ is $\sin\varphi$, so its circumference is
	$2\pi\sin\varphi$. Drawing the cone tangent to $S^2$ at $C$,
	the length from the apex of the cone to a point on $C$ is $\tan\varphi$.
	If we ``unroll'' the cone into a sector of a circle
	with radius $\tan\varphi$ by cutting it along one of its generators,
	we can compute that the angle of the sector is $2\pi\cos\varphi$
	by using the arc length and radius values;
	this is what gives us $\Delta\varphi = 2\pi - 2\pi\cos\theta$.
	
	For the second part of the question we can calculate $A$
	with spherical coordinates to get
	\[A = \int_0^{2\pi}\int_0^\varphi \sin\phi\dif\psi\dif\theta
	= 2\pi(1 - \cos\theta)\]
	which gives the result...?
\end{soln}


















\end{document}